\documentclass[11pt,a4paper]{article}
\usepackage[T1]{fontenc}
\usepackage{geometry}
\usepackage{amsmath}
\usepackage{amssymb}
\usepackage{booktabs}
\usepackage{graphicx}
\usepackage{xcolor}
\usepackage{array}
\usepackage{mathptmx}
\usepackage{listings}
\usepackage{hyperref}

\geometry{margin=1in}

\definecolor{codebg}{gray}{0.95}
\lstset{
  basicstyle=\ttfamily\small,
  backgroundcolor=\color{codebg},
  breaklines=true,
  frame=single,
  columns=fullflexible,
  showstringspaces=false
}

\title{Running a Reaction-Diffusion GPU Solver under Nixnan\\
\large FP64 vs FP32 Dynamic Range, Exp-Binade Histogramming, and an
Exponential \texttt{SAMPLING} Sweep}
\author{SC26 Tutorial / Nixnan Experiments}
\date{}

\begin{document}
\maketitle

\section{Executive Summary}

This report documents \texttt{reaction\_diffusion\_gpu}, a CUDA port of the
1D reaction-diffusion solver from the FPChecker tutorial
(\texttt{tutorial/example\_2/reaction\_diffusion.cpp} in the sibling
\texttt{SC26-Tutorial-NixNan} repository), run under nixnan
(\url{https://github.com/parfloat/nixnan}) on an NVIDIA RTX 3090. It
replaces the repository's prior ad-hoc \texttt{rd\_nixnan.cu} example
(which used its own invented FP16/BF16/FP32 comparison and parameters) with
a solver that follows the tutorial's PDE and parameters exactly, and reuses
the exp-binade histogram / exponential \texttt{SAMPLING} experiment pattern
established for \texttt{lu\_solve\_gpu}
(\texttt{../lu\_solve\_gpu/lu\_solve\_gpu\_report.tex}). Key results:

\begin{enumerate}
  \item \textbf{FP64 vs FP32 dynamic range}: with $\lambda=25$ amplifying
    the solution as $\sim e^{\lambda t}$, the growth factor by $t=4$ is
    $e^{100}\approx 2.7\times10^{43}$ -- inside FP64's range but well past
    FP32's ($\approx 3.4\times10^{38}$). The FP64 build completes cleanly
    (max $1.70\times10^{43}$ at $t=4$, zero exceptions); the FP32 build
    overflows to Infinity between $t=3.52$ and $t=3.60$, and the resulting
    NaN then propagates for the rest of the run -- reproducing the CPU
    tutorial's central lesson (``this case should remain in FP64'') on the
    GPU.
  \item \textbf{A silent-failure wrinkle}: the GPU port's max-reduction
    mirrors the CPU version's \texttt{std::max\_element}, including its
    NaN-comparison hazard. In this program that hazard goes one step
    further than the CPU tutorial's own commentary anticipates: it does
    not just report a misleading value, it makes the reported max
    \emph{look finite} (literally \texttt{0.0}) from the moment of overflow
    onward, silently defeating the program's own \texttt{std::isfinite()}
    check. nixnan, which instruments individual FADD/FFMA instructions and
    never looks at the reduction's output, catches the real Infinity and
    NaN regardless (Section~\ref{sec:silent-failure}).
  \item \textbf{Same tool-level findings as \texttt{lu\_solve\_gpu}
    reproduce here}: the \texttt{SAMPLING} sweep again shows byte-identical
    logs across levels $1,2,4,8$, consistent with the key-mismatch bug
    documented in \texttt{../lu\_solve\_gpu/lu\_solve\_gpu\_report.tex}
    (Section~\ref{sec:sampling}).
  \item \textbf{Two additional nixnan features}, not exercised in the
    \texttt{lu\_solve\_gpu} experiment, are demonstrated here:
    \texttt{PRINT\_ILL\_INSTR}+\texttt{MAX\_ERRORS} source/SASS-attributed
    diagnostics, and \texttt{INSTR\_MEM} memory instrumentation
    (Section~\ref{sec:stage3}).
\end{enumerate}

\section{Provenance and Algorithm}

\texttt{reaction\_diffusion\_gpu.cu} solves the 1D linear reaction-diffusion
equation
\[
\frac{\partial u}{\partial t} = D\,\frac{\partial^2 u}{\partial x^2} + \lambda u,
\qquad u(0,t)=u(L,t)=0, \qquad u(x,0)=\sin\!\left(\frac{\pi x}{L}\right),
\]
with an explicit finite-difference (FTCS) update, using exactly the
parameters from \texttt{tutorial/example\_2/reaction\_diffusion.cpp}:
$L=1$, $T=4$, $N=101$ spatial points, $D=0.01$, $\lambda=25$, $M=80000$
time steps ($dt = T/M = 5\times10^{-5}$). One CUDA kernel
(\texttt{rd\_step\_kernel}) performs the update for all interior points in
parallel each step; a second kernel (\texttt{max\_reduce\_kernel}) computes
the running maximum, launched every $M/50 = 1600$ steps to match the CPU
version's print cadence.

Precision is a compile-time macro (\texttt{RD\_USE\_FLOAT}) rather than the
CPU tutorial's hand-edited \texttt{typedef}, so both variants exist as
separate binaries (\texttt{reaction\_diffusion\_gpu\_fp64},
\texttt{reaction\_diffusion\_gpu\_fp32}) built from the same source by the
\texttt{Makefile} -- this is what \texttt{run\_nixnan\_experiment.sh}
automates instead of requiring a manual edit-and-rebuild cycle per
precision.

\subsection{Build and Environment}
\begin{itemize}
  \item \textbf{GPU}: NVIDIA GeForce RTX 3090, compute capability 8.6
  \item \textbf{Compilation}: \texttt{nvcc -std=c++17 -g -O2 -lineinfo -arch=sm\_86}
  \item \textbf{Instrumentation}: \texttt{LD\_PRELOAD=../../nixnan.so}
\end{itemize}

\section{FP64 vs FP32: Dynamic Range}
\label{sec:dynamic-range}

\begin{table}[ht]
\centering
\begin{tabular}{lrr}
\toprule
& FP64 & FP32 \\
\midrule
Max $|u|$ at $t=3.52$ & $1.11\times10^{38}$ & $1.10\times10^{38}$ \\
Max $|u|$ at $t=3.60$ & $8.09\times10^{38}$ & \texttt{0.0} (see \S\ref{sec:silent-failure}) \\
Max $|u|$ at $t=4.00$ & $1.70\times10^{43}$ & \texttt{0.0} \\
First non-finite reported by the program itself & never & never (masked) \\
FP32/FP64 exceptions (nixnan, Stage 1) & 0/0/0 & 36 NaN / 16 Inf / 40 $-$Inf \\
\bottomrule
\end{tabular}
\caption{From \texttt{logs/baseline/\{fp64,fp32\}.log} and
\texttt{logs/histogram/\{fp64,fp32\}.log}. FP32's reported maxima are the
reduction's output, not the true array contents -- see
Section~\ref{sec:silent-failure}.}
\label{tab:dynamic-range}
\end{table}

Both builds agree closely up through $t=3.52$ (max $\approx1.1\times10^{38}$,
just below FP32's ceiling of $\approx3.4\times10^{38}$). By the next printed
sample, $t=3.60$, FP32 has overflowed to Infinity while FP64 continues to
grow smoothly; FP64 finishes the full simulation at $1.70\times10^{43}$,
comfortably inside its $\approx1.8\times10^{308}$ range. This is exactly the
tutorial's intended lesson, now demonstrated on the GPU: a workload that
"looks fine" while values stay under $\sim10^{38}$ becomes numerically
unsafe in FP32 the moment it doesn't.

\section{A Silent-Failure Wrinkle in the GPU Port}
\label{sec:silent-failure}

The CPU version's \texttt{print\_max\_value} uses \texttt{std::max\_element},
whose pairwise comparisons are all \texttt{false} once a NaN is involved (a
NaN is neither greater than nor less than anything, including itself).
\texttt{max\_reduce\_kernel} in this port intentionally mirrors that logic
(\texttt{v > best}) rather than "fixing" it with e.g.\ \texttt{fmax}, in
order to carry the same authentic hazard onto the GPU.

In this program, however, the hazard manifests more severely than a merely
misleading value. \texttt{logs/baseline/fp32.log} shows the reported
maximum becoming exactly

\begin{lstlisting}
At time t = 3.600000, Maximum value of u = 0.0000000000e+00
At time t = 3.680000, Maximum value of u = 0.0000000000e+00
...
At time t = 4.000000, Maximum value of u = 0.0000000000e+00
\end{lstlisting}

\noindent from $t=3.60$ onward -- not \texttt{nan}, not \texttt{inf}, but a
literal, plausible-looking, finite \texttt{0.0}. The mechanism: the
reduction seeds \texttt{best} from \texttt{u[0]} (a boundary point, pinned
to exactly \texttt{0.0} by the Dirichlet condition) and only updates it via
\texttt{v > best}. Once the array's interior values become
Infinity/NaN, none of those comparisons succeed, so \texttt{best} never
leaves its seed value of \texttt{0.0}.

This has a second-order consequence: \texttt{reaction\_diffusion\_gpu.cu}
also contains a self-check --

\begin{lstlisting}[language=C++]
if (!reported_nonfinite && !std::isfinite((double)max_val))
{
    reported_nonfinite = true;
    std::cerr << "first non-finite reported max at step " << ...
}
\end{lstlisting}

\noindent This message never appears in \texttt{logs/baseline/fp32.log}'s
stderr, because \texttt{max\_val} itself is \texttt{0.0} -- perfectly
finite by the check's own definition -- even though the array it was
reduced from genuinely contains Infinity and NaN. A monitoring strategy
built entirely on top of this program's own reported summary statistic
would conclude, incorrectly, that nothing ever went wrong.

nixnan is unaffected by this because it does not consume the reduction's
output at all: it instruments every \texttt{FADD}/\texttt{FFMA} as it
executes, inside \texttt{rd\_step\_kernel} itself, before the corrupted
value ever reaches \texttt{max\_reduce\_kernel}. \texttt{logs/histogram/fp32.log}
and \texttt{logs/memory/fp32.log} both report the real counts (36 NaN, 16
Infinity, 40 $-$Infinity in registers; 1 NaN value, 36{,}928 repeats,
actually written to global memory) that the program's own console output
never shows. This is, in miniature, the argument for binary-level
floating-point exception instrumentation: a program's own summary
statistics can be fooled by exactly the kind of comparison logic that
silently ignores non-finite values, while instruction-level instrumentation
cannot.

\section{Stage 1: Exp-Binade Histogram Gathering}
\label{sec:histogram}

\texttt{bin\_spec.json} (generated by \texttt{gen\_bin\_spec.py}) populates
\emph{both} \texttt{f32} and \texttt{f64} -- unlike \texttt{lu\_solve\_gpu},
where only one precision existed per build, here both formats are
genuinely present across the two binaries that share this one spec file.
Each format's full valid exponent range is tiled in buckets of width 4
(the same "bucket size 4" convention as
\texttt{../lu\_solve\_gpu/gen\_bin\_spec.py}): \texttt{f32}
$\in[-126,127]$ (64 buckets), \texttt{f64}$\in[-1022,1023]$ (512 buckets).
\texttt{f16}/\texttt{bf16} are left empty; this program never uses them.

The report-threshold \texttt{count} required re-tuning for this workload.
$M=80000$ steps over $\sim$99 interior points produce roughly $1000\times$
more instrumented FP-op instances than \texttt{lu\_solve\_gpu}'s one-shot
$20\times20$ solve; reusing \texttt{lu\_solve\_gpu}'s \texttt{count: 128}
produced $\sim$950,000 report lines here (Table~\ref{tab:count-tuning}).
\texttt{count: 1{,}000{,}000} was chosen empirically to land back in the
same "nice" $\sim$100--150-reports-per-run range.

\begin{table}[ht]
\centering
\begin{tabular}{rrr}
\toprule
\texttt{count} & FP64 reports & FP32 reports \\
\midrule
128 & 1{,}050{,}603 & 943{,}683 \\
100{,}000 & 1{,}339 & 1{,}203 \\
1{,}000{,}000 (chosen) & 121 & 109 \\
\bottomrule
\end{tabular}
\caption{Report-line volume for a single run at three candidate
\texttt{count} thresholds (measured during tuning; not stored under
\texttt{logs/}).}
\label{tab:count-tuning}
\end{table}

\section{Stage 2: Exponential \texttt{SAMPLING} Sweep}
\label{sec:sampling}

The same schedule as \texttt{../lu\_solve\_gpu/run\_nixnan\_experiment.sh}
was applied to \texttt{reaction\_diffusion\_gpu\_fp32}: \texttt{SAMPLING}
$\in\{1,2,4,8\}$, 16 process launches ("snaps") per level, then reset back
to \texttt{1} -- 64 whole-program runs, stored under
\texttt{logs/sampling/cycle\_01/sampling\_\{1,2,4,8\}/snap\_\{01..16\}}.
\texttt{rd\_step\_kernel} is launched 80{,}000 times per run, so
\texttt{SAMPLING} has far more opportunity here than in
\texttt{lu\_solve\_gpu} (19 launches per kernel) to visibly change
instrumentation coverage if it worked. It still does not:
\texttt{sampling\_1/snap\_01.log} and \texttt{sampling\_8/snap\_01.log} are
byte-identical, consistent with the key-mismatch bug documented in
\texttt{../lu\_solve\_gpu/lu\_solve\_gpu\_report.tex}
(Section~4.2 there: \texttt{should\_instrument()} looks up the per-kernel
invocation counter by the full mangled signature, while
\texttt{nvbit\_at\_cuda\_event()} increments it under the name truncated by
\texttt{cut\_kernel\_name()}, so the two keys never match). This report does
not re-derive that finding in full; see the earlier report for the code
excerpts.

\section{Stage 3: Interesting Exception Settings}
\label{sec:stage3}

Two nixnan features not exercised in the \texttt{lu\_solve\_gpu} experiment:

\subsection{3a: \texttt{PRINT\_ILL\_INSTR} + \texttt{MAX\_ERRORS}}

Run with \texttt{PRINT\_ILL\_INSTR=1 MAX\_ERRORS=20}, nixnan prints the
exact offending SASS instruction for each exception and terminates the
process once 20 have been reported. On \texttt{reaction\_diffusion\_gpu\_fp32}
this cap is reached almost immediately after the overflow at $t\approx3.6$;
of the 20 reported errors, 14 are attributed to
\texttt{reaction\_diffusion\_gpu.cu:64} (the diffusion-term computation,
\texttt{D\_dt\_over\_dx2 * (u[i+1] - 2*u[i] + u[i-1])}) and 6 to
\texttt{reaction\_diffusion\_gpu.cu:66} (the state update,
\texttt{u\_next[i] = u[i] + diffusion\_term + reaction\_term}) --
consistent with $\text{Inf} - 2\cdot\text{Inf} + \text{Inf} = \text{NaN}$
occurring right where the stencil combines three already-overflowed
neighbors. \texttt{reaction\_diffusion\_gpu\_fp64} never reaches the cap
and runs to completion (\texttt{logs/diagnostic/fp64.log} has zero error
lines).

\subsection{3b: \texttt{INSTR\_MEM}}

Run with \texttt{INSTR\_MEM=1}, nixnan additionally instruments memory
store instructions. \texttt{logs/memory/fp32.log} shows

\begin{lstlisting}
#nixnan: --- FP32 Memory  Operations ---
#nixnan: NaN:                    1 (36928 repeats)
\end{lstlisting}

\noindent confirming that NaN is not merely a transient register value but
is actually written to \texttt{u\_next} in GPU global memory, repeatedly,
for the remainder of the run (80{,}000 total steps, overflow first appears
around step 72{,}000 at $t\approx3.6$, so on the order of $8000\times$
99 $\approx 800{,}000$ store opportunities remain -- 36{,}928 of the
distinct memory NaN events were captured). \texttt{logs/memory/fp64.log}
shows zero, as expected.

\section{Conclusions}

\begin{enumerate}
  \item The GPU port reproduces the CPU tutorial's central lesson exactly:
    FP64 is safe for this workload's full time horizon; FP32 is not, and
    fails via the textbook overflow-then-NaN-propagation path.
  \item The port's naive max-reduction (chosen deliberately to mirror
    \texttt{std::max\_element}'s NaN hazard) turns out to hide the FP32
    overflow from the program's own console output and its own
    finiteness self-check -- a instructive illustration of why
    instruction-level exception detection is valuable independent of
    whatever a program's own diagnostics report.
  \item The \texttt{SAMPLING} key-mismatch bug found while working on
    \texttt{lu\_solve\_gpu} reproduces here under a workload with $1000\times$
    more kernel launches, reinforcing that it is a property of nixnan
    itself and not an artifact of the smaller LU experiment.
  \item \texttt{PRINT\_ILL\_INSTR}+\texttt{MAX\_ERRORS} and
    \texttt{INSTR\_MEM} both worked as documented and produced exactly the
    kind of source-line- and memory-level evidence the CPU tutorial's own
    \texttt{fpc-create-report} is meant to provide, here on the GPU build.
\end{enumerate}

\section*{Reproduction}

\begin{lstlisting}
cd ~/repos/nixnan/SC26-expts/rd_nixnan
./run_nixnan_experiment.sh
\end{lstlisting}

\noindent Requires a built \texttt{../../nixnan.so}. See \texttt{README.md}
in this directory for file-by-file documentation and override variables
(\texttt{NIXNAN\_SO}, \texttt{CYCLES}, \texttt{SAMPLING\_EXE}).

\end{document}
